\documentclass[12pt]{article}

% --- Packages ---
\usepackage[utf8]{inputenc}
\usepackage[T1]{fontenc}
\usepackage[french]{babel}
\usepackage{amsmath, amssymb}
\usepackage{geometry}
% \usepackage[draft]{graphicx}
\usepackage{tikz}
\usepackage[most]{tcolorbox}
\usepackage{multicol}
\usepackage{lmodern}
\usepackage{pgfplots}
\pgfplotsset{compat=1.18}
\usepackage{mdframed}
\usepackage{array}
\usepackage[thinlines]{easytable}
\usepackage{tkz-tab}
\usepackage{bm}
\usepackage{xcolor}
\usepackage{sectsty}
\usepackage[inline]{enumitem}
\usepackage[dvipsnames]{xcolor}
\usepackage{multido}
% --- Page ---
\geometry{margin=2cm}

% Footer
\usepackage{fancyhdr}
\pagestyle{fancy}
\cfoot{\thepage}
\rfoot{\textit{G.M.}}
\renewcommand{\headrulewidth}{0pt}
\renewcommand{\footrulewidth}{0pt}

\setlength{\parindent}{0pt}

% Les cadres ne sont jamais coupés entre deux pages
\mdfsetup{nobreak=true}

\sectionfont{\color{red}\Large\bfseries}     % pour \section*
\subsectionfont{\color{red}\large\bfseries\hspace{1cm}}

\begin{document}

% ============================
%   Calcul littéral
% ============================

\begin{center}
    {\LARGE \textbf{TD - Calcul littéral}}\\[4mm]
\end{center}

% ============================
%   Partie 1 : Développer
% ============================

\section*{\textcolor{red}{I/ Développer une expression}}

\vspace{5pt}
\begin{mdframed}[
    linewidth=1pt,
    roundcorner=5pt,
    linecolor=Red,
    innertopmargin=10pt,
    innerbottommargin=10pt
]
{\color{red}
\textbf{\underline{Propriété (simple distributivité) :}}

\vspace{5pt}
Soient $k$, $a$ et $b$ trois nombres. Alors :

\vspace{10pt}
\begin{center}
    $\bm{k(a+b) = }$ \_\_\_\_\_\_\_\_\_\_\_\_ \hspace{2cm} $\bm{k(a-b) = }$ \_\_\_\_\_\_\_\_\_\_\_\_
\end{center}
}
\end{mdframed}

\vspace{5pt}
\underline{Exemple :} $3(x+5) =$ \_\_\_\_\_\_\_\_\_\_\_\_\_\_\_\_\_\_\_\_\_\_\_\_

\vspace{10pt}
\underline{Remarque :} Un signe $-$ devant une parenthèse change \textbf{tous} les signes à l'intérieur.

\vspace{10pt}
\hspace{1cm}$-(a+b) = $ \_\_\_\_\_\_\_\_\_\_\_\_

\vspace{10pt}
\begin{mdframed}[
    linewidth=1pt,
    roundcorner=5pt,
    linecolor=RoyalBlue,
    innertopmargin=10pt,
    innerbottommargin=10pt
]
{\color{RoyalBlue}
\textbf{\underline{Exercice d'application n°1:} Développer avec la simple distributivité.}}

\vspace{5pt}
Développer puis réduire les expressions suivantes :

\vspace{5pt}
\hspace{1cm}a) $A = 4(x+2)$ \hspace{3cm}b) $B = -3(2x-5)$

\vspace{0.5cm}
{\color{black!45}\dotfill}

\vspace{0.5cm}
{\color{black!45}\dotfill}

\vspace{0.5cm}
{\color{black!45}\dotfill}

\vspace{0.5cm}
{\color{black!45}\dotfill}

\vspace{20pt}

\end{mdframed}

\begin{mdframed}[
    linewidth=1pt,
    roundcorner=5pt,
    linecolor=Red,
    innertopmargin=10pt,
    innerbottommargin=10pt
]
{\color{red}
\textbf{\underline{Propriété (double distributivité) :}}

\vspace{5pt}
Soient $a$, $b$, $c$ et $d$ quatre nombres. Alors :

\vspace{10pt}
\begin{center}
    $\bm{(a+b)(c+d) = }$ \_\_\_\_\_\_\_\_\_\_\_\_\_\_\_\_\_\_\_\_\_\_\_\_
\end{center}
}
\end{mdframed}

\vspace{5pt}
\underline{Exemple :} $(x+2)(x+3) = $ \_\_\_\_\_\_\_\_\_\_\_\_\_\_\_\_\_\_\_\_\_\_\_\_ $=$ \_\_\_\_\_\_\_\_\_\_\_\_

\vspace{10pt}
\begin{mdframed}[
    linewidth=1pt,
    roundcorner=5pt,
    linecolor=RoyalBlue,
    innertopmargin=10pt,
    innerbottommargin=10pt
]
{\color{RoyalBlue}
\textbf{\underline{Exercice d'application n°2:} Développer avec la double distributivité.}}

\vspace{5pt}
Développer puis réduire les expressions suivantes :

\vspace{5pt}
\hspace{1cm}a) $E = (x+3)(x+5)$ \hspace{2cm}b) $F = (3-x)(2x+1)$

\vspace{0.5cm}
{\color{black!45}\dotfill}

\vspace{0.5cm}
{\color{black!45}\dotfill}

\vspace{0.5cm}
{\color{black!45}\dotfill}

\vspace{0.5cm}
{\color{black!45}\dotfill}

\vspace{0.5cm}
{\color{black!45}\dotfill}

\vspace{20pt}

\end{mdframed}

\vspace{5pt}
% ============================
%   Partie 2 : Factoriser
% ============================
\section*{\textcolor{red}{II/ Factoriser une expression}}

\vspace{5pt}
\underline{Remarque :} Factoriser, c'est l'opération \textbf{inverse} de {\color{black!45}\dotfill}

\vspace{10pt}
\begin{mdframed}[
    linewidth=1pt,
    roundcorner=5pt,
    linecolor=Red,
    innertopmargin=10pt,
    innerbottommargin=10pt
]
{\color{red}
\textbf{\underline{Propriété :}}

\vspace{5pt}
Soient $k$, $a$ et $b$ trois nombres. Le nombre $k$ est appelé \_\_\_\_\_\_\_\_\_\_\_\_\_\_\_\_.

\vspace{5pt}
\begin{center}
    $\bm{ka + kb = }$ \_\_\_\_\_\_\_\_\_\_\_\_ \hspace{2cm} $\bm{ka - kb = }$ \_\_\_\_\_\_\_\_\_\_\_\_
\end{center}
}
\end{mdframed}

\vspace{5pt}
\underline{Exemples :}

\vspace{5pt}
\textbullet \ $5x + 15 = $ \_\_\_\_\_\_\_\_\_\_\_\_\_\_\_\_\_\_\_\_\_\_\_\_\_\_\_\_\_\_\_\_\_\_\_\_

\vspace{5pt}
\begin{mdframed}[
    linewidth=1pt,
    roundcorner=5pt,
    linecolor=RoyalBlue,
    innertopmargin=10pt,
    innerbottommargin=10pt
]
{\color{RoyalBlue}
\textbf{\underline{Exercice d'application n°3:} Factoriser avec un facteur commun simple.}}

\vspace{5pt}
Factoriser les expressions suivantes :

\vspace{5pt}
\hspace{1cm}a) $A = 10x - 15$ \hspace{3cm}b) $B = 8x^2 - 12x$

\vspace{0.5cm}
{\color{black!45}\dotfill}

\vspace{0.5cm}
{\color{black!45}\dotfill}

\vspace{0.5cm}
{\color{black!45}\dotfill}

\vspace{20pt}

\end{mdframed}

\vspace{5pt}
\textbf{Le facteur commun peut aussi être une parenthèse :}

\vspace{5pt}
\underline{Exemple :}

\vspace{-10pt}
\begin{align*}
    (x+1)(2x+3) + (x+1)(x-5) &= (x+1)\big[\ \dots\dots\dots\dots\ + \ \dots\dots\dots\dots\ \big] \\[6pt]
    &= (x+1)(\dots\dots\dots\dots)
\end{align*}

\underline{Remarque :} Attention au facteur « caché » : $(2x+1) = (2x+1) \times $ \_\_\_\_\_\_

\vspace{10pt}
\begin{mdframed}[
    linewidth=1pt,
    roundcorner=5pt,
    linecolor=RoyalBlue,
    innertopmargin=10pt,
    innerbottommargin=10pt
]
{\color{RoyalBlue}
\textbf{\underline{Exercice d'application n°4:} Factoriser avec une parenthèse en facteur commun.}}

\vspace{5pt}
Factoriser puis réduire les expressions suivantes :

\vspace{5pt}
\hspace{1cm}b) $A = (2x-3)(x+1) - (2x-3)(4x+5)$

\vspace{5pt}
\hspace{1cm}c) $B = (x-4)^2 + (x-4)(2x+1)$

\vspace{5pt}
\hspace{1cm}d) $C = (3x+1) + (3x+1)(x-2)$

\vspace{0.5cm}
{\color{black!45}\dotfill}

\vspace{0.5cm}
{\color{black!45}\dotfill}

\vspace{0.5cm}
{\color{black!45}\dotfill}

\vspace{0.5cm}
{\color{black!45}\dotfill}

\vspace{0.5cm}
{\color{black!45}\dotfill}

\vspace{0.5cm}
{\color{black!45}\dotfill}

\vspace{0.5cm}
{\color{black!45}\dotfill}

\vspace{0.5cm}
{\color{black!45}\dotfill}

\vspace{20pt}

\end{mdframed}

% ============================
%   Partie 3 : Identités remarquables
% ============================

\section*{\textcolor{red}{III/ Identités remarquables}}

\begin{mdframed}[
    linewidth=1pt,
    roundcorner=5pt,
    linecolor=Red,
    innertopmargin=10pt,
    innerbottommargin=10pt
]
{\color{red}
\textbf{\underline{Propriété :}}

\vspace{5pt}
Soient $a$ et $b$ deux nombres. Alors :

\vspace{-20pt}
\begin{center}
    \vspace{20pt}
    $\bm{(a+b)^2 = }$ \_\_\_\_\_\_\_\_\_\_\_\_\_\_\_\_\_\_\_\_

    \vspace{20pt}
    $\bm{(a-b)^2 = }$ \_\_\_\_\_\_\_\_\_\_\_\_\_\_\_\_\_\_\_\_

    \vspace{20pt}
    $\bm{(a-b)(a+b) = }$ \_\_\_\_\_\_\_\_\_\_\_\_\_\_\_\_\_\_\_\_
\end{center}
}
\end{mdframed}

\begin{mdframed}[
    linewidth=1pt,
    roundcorner=5pt,
    linecolor=Plum,
    innertopmargin=10pt,
    innerbottommargin=10pt
]
{\color{Plum}
\textbf{\underline{Démonstration :}}}

\vspace{5pt}
Démontrons la première égalité à l'aide de la double distributivité.

\vspace{0.5cm}
$(a+b)^2 = ${\color{black!45}\dotfill}

\vspace{0.5cm}
{\color{black!45}\dotfill}

\vspace{10pt}

\end{mdframed}

\vspace{5pt}
\underline{Exemple :} $(x+3)^2 = $ \_\_\_\_\_\_\_\_\_\_\_\_\_\_\_\_\_\_\_\_\_\_\_\_\_\_\_\_\_\_\_\_\_\_\_\_

\begin{mdframed}[
    linewidth=1pt,
    roundcorner=5pt,
    linecolor=RoyalBlue,
    innertopmargin=10pt,
    innerbottommargin=10pt
]
{\color{RoyalBlue}
\textbf{\underline{Exercice d'application n°5:} Développer avec les identités remarquables.}}

\vspace{5pt}
Développer puis réduire les expressions suivantes :

\vspace{5pt}
\hspace{1cm}a) $A = (x+4)^2$ \hspace{2cm}b) $B = (3x-2)^2$ \hspace{2cm} $T = (x+7)(x-7)$

\vspace{0.5cm}
{\color{black!45}\dotfill}

\vspace{0.5cm}
{\color{black!45}\dotfill}

\vspace{0.5cm}
{\color{black!45}\dotfill}

\vspace{0.5cm}
{\color{black!45}\dotfill}

\vspace{20pt}

\end{mdframed}

\vspace{5pt}
\textbf{Les identités remarquables se lisent aussi de droite à gauche : elles permettent alors }

\vspace{10pt}
\textbf{de \_\_\_\_\_\_\_\_\_\_\_\_.}

\vspace{5pt}
\underline{Exemples :}

\vspace{5pt}
\textbullet \ $x^2 - 6x + 9 =$ \_\_\_\_\_\_\_\_\_\_\_\_\_\_\_\_\_\_\_\_\_\_\_\_\_\_\_\_\_\_\_\_\_\_\_\_

\begin{mdframed}[
    linewidth=1pt,
    roundcorner=5pt,
    linecolor=RoyalBlue,
    innertopmargin=10pt,
    innerbottommargin=10pt
]
{\color{RoyalBlue}
\textbf{\underline{Exercice d'application n°6:} Factoriser avec les identités remarquables.}}

\vspace{5pt}
Factoriser les expressions suivantes :

\vspace{5pt}
\hspace{1cm}a) $A = 9x^2 - 16$ \hspace{1cm}b) $B = x^2 + 10x + 25$ \hspace{1cm}e) \textit{Pour aller plus loin :} $C = (x+1)^2 - 9$

\vspace{0.5cm}
{\color{black!45}\dotfill}

\vspace{0.5cm}
{\color{black!45}\dotfill}

\vspace{0.5cm}
{\color{black!45}\dotfill}

\vspace{0.5cm}
{\color{black!45}\dotfill}

\vspace{0.5cm}
{\color{black!45}\dotfill}

\vspace{10pt}

\end{mdframed}

\end{document}
